\documentclass[book.tex]{subfiles}
\begin{document}

\chapter{Implicit Methods and Stiffness in Ordinary Differential Equations}

\label{chap:implicit_methods}
\noindent\rule{11cm}{0.4pt} \\
\textbf{Prerequisites:} \autoref{chap:calculus} \\ \noindent\rule{11cm}{0.4pt}
\vspace{5mm}

In this chapter we will learn what stiffness means and the implications stiffness has for solving ordinary differential equations. We will then learn how the implicit Euler's method can be used to solve stiff ordinary differential equations.

\section{Stiff ODEs}
Stiffness is a subtle, difficult, and important concept in the numerical solution of ordinary differential equations. It depends on the differential equation, the initial conditions, and the numerical method. A \textit{stiff} system is one which involves rapidly changing components together with slowly changing ones. In some cases, the rapidly varying components are ephemeral (lasting a short time) transients that die away quickly, after which the solution becomes dominated by the slowly varying components. Although the transient phenomena exist for only a short part of the integration interval, they can dictate the time step for the entire solution.

Both individual ODEs and system of ODEs can be stiff. A simple example of a stiff ODE is given below for visualization.
\begin{align}
	\frac{dy}{dt} = -1000y+3000-2000e^{-t}
\end{align}
Let the initial condition be $y(0)=0$. The analytical solution for this problem can be solved as 
\begin{align}
y(t) = 3 - 0.998e^{-1000t}-2.002e^{-t}
\end{align}
\begin{marginfigure}[5\baselineskip] 
	\centering
	% \includegraphics[width = 2.5in]{figures/part1b/implicit_methods/stiff.PNG}
 \includegraphics[width = 2.5in]{figures/part1b/implicit_methods/stiff_0.png}
	\centering
	\caption{Analytical solution of $	\frac{dy}{dt} = -1000y+3000-2000e^{-t}$}
	\label{fig:1}
\end{marginfigure}

This solution is plotted and as it can be seen, the solution is initially dominated by the fast exponential term of $e^{-1000t}$. Although the solution appears to start at 1, there is actually a fast transient from $y$ = $0$ to $1$ that occurs in less than the 0.005 time unit. This transient is perceptible only when the response is viewed on the finer timescale in the inset. After a short period (t < 0.005), this transient dies out and the solution becomes governed by the slow exponential ($e^{-t}$). 

Insight into the step size required for stability of such a solution can be gained by examining the homogeneous part of the differential equation.
\begin{align}
\frac{dy}{dt} = -1000y
\end{align}
If $y(0)=y_0$, calculus can be used to determine the solution as $y = y_0e^{-1000t}$. Thus, the solution starts at $y_0$ and asymptotically approaches zero. Euler's method can be used to solve the same problem numerically:
\begin{align}
y_{i+1} = y_i + \dfrac{dy_i}{dt}h
\end{align}
For this ODE, the above equation is simplified to
\begin{align}
y_{i+1} = y_i -1000y_ih = y_i(1-1000h)
\end{align}
The stability of this formula clearly depends on the step size $h$.	That is, $|1-1000h|$ must be less than $1$. Thus, if $h>2/1000$, $|y_i| \rightarrow \infty$ as $i \rightarrow \infty$. For the fast transient part of the solution ($e^{-1000t}$ term), this criterion shows that the step size to maintain stability must be less than $2/1000 = 0.002$. In addition, it should be noted that, whereas this criterion maintains stability (i.e., a bounded solution), an even smaller step size would be required to obtain an accurate solution. Thus, although the transient occurs for only a small fraction of the integration interval, it controls the maximum allowable step size.


\section{Implicit methods and Adaptive time stepping}

One way of solving the above problem is to use adaptive time-stepping, to avoid the small time-step throughout. This is obtained by bounding the local truncation error as discussed in earlier chapter. Some Physika code for finding the adaptive time steps to solve the integral is given below.

\begin{marginfigure}
	\centering
	\includegraphics[width = 2.5in]{figures/part1b/implicit_methods/adaptive.png}
	\centering
	\caption{Step-Size Plot for the Adaptive Euler's Method}
	\label{fig:5}
\end{marginfigure}


\begin{lstlisting}[language=python]
def adaptive_step($\frac{dy}{dt}$: $\mathbb{R} \to \mathbb{R}$, t$_i$: $\mathbb{R}$ = 0, t$_f$: $\mathbb{R}$ = 3, t_adp: $\mathbb{R}[n]$):
  f_d = $\nabla$($\frac{dy}{dt}$, t)
  y[0] = 0, e_b = 1e-10, t_adp[0] = t$_i$
  h = sqrt(2*e_b/|f_d(y[0],t_adp[0])|)
  while(t_adp[i] < t$_f$):
    y[i+1] = y[i] + $\frac{dy}{dt}$(y[i],t_adp[i]) * h[i]
    h[i+1] = sqrt(2*e_b/abs(f_d(y[i],t_adp[i])))
    t_adp[i+1] = t_adp[i]+h[i+1]
    i = i+1
\end{lstlisting}

In contrast to explicit approaches, implicit methods offer an alternative remedy.
Such representations are called implicit because the unknown appears on both sides of the equation. An implicit form of Euler’s method can be developed by evaluating the derivative at the future time. An implicit form of Euler's method would be:
\begin{align}
y_{i+1} = y_i + \dfrac{dy_{i+1}}{dt}h
\end{align}


This is called the \textit{backward} or \textit{implicit} Euler's method.
For the ODE under consideration, the above equation is simplified to
\begin{align}
y_{i+1} = y_i - 1000y_{i+1}h
\end{align}
which can be solved for as
\begin{align}
y_{i+1} = \dfrac{y_i}{1+1000h}
\end{align}
For this case, regardless of the size of the step, $|y_i|\rightarrow 0$ as $i\rightarrow \infty$. Hence, the approach is called \textit{unconditionally stable}.

\begin{marginfigure}[\baselineskip] 
	\centering
	% \includegraphics[width = 2.5in]{figures/part1b/implicit_methods/stiffeuler.png}
 \includegraphics[width = 2.5in]{figures/part1b/implicit_methods/stiffeuler0.png}
	\centering
	\caption{Solution of a stiff ODE with (a) the explicit (without time-stepping) and (b) implicit Euler methods}
	\label{fig:2}
\end{marginfigure}

For linear ODEs, the equations need to be rearranged to solve for the next ODE update in the implicit Euler's method approach. For nonlinear ODEs, the solution becomes even more difficult since it involves solving a system of nonlinear simultaneous equations. Even though stability is gained through implicit approaches, it is at the cost of added solution complexity. We will look at how to solve such a non-linear system in the example below.

\section{Van der Pol Equation}
The van der Pol equation is a model of an electronic circuit that
arose back in the days of vacuum tubes. It evolves in time according to the second-order differential equation:
\begin{align}
\dfrac{d^2y}{dt^2} - \mu(1-y^2)\dfrac{dy}{dt} + y = 0
\end{align}

The solution to this equation becomes progressively stiffer as $\mu$ gets large. As an example, given the initial conditions, $y(0) = 0$ and $\dfrac{dy}{dt} = 1$ at $t=0$, we want to solve the van der Pol equation using explicit and implicit Euler's Method from $t=0$ to $t=10$, using a step-size of 0.05 and $\mu$ of 5. 

The first step is to convert the second-order ODE into a pair of first-order ODEs by defining
\begin{marginfigure}
	\centering
	\includegraphics[width = 2.5in]{figures/part1b/implicit_methods/explicit.png}
	\centering
	\caption{Since the van der Pol equation with $\mu$ = 5 is a stiff system, the explicit method clearly fails in capturing the response of the variable.}
	\label{fig:3}
\end{marginfigure}

\begin{marginfigure}
	\centering
	\includegraphics[width = 2.5in]{figures/part1b/implicit_methods/implicit.png}
	\centering
	\caption{Solution of the van der Pol equation using the Euler's Implicit Method}
	\label{fig:4}
\end{marginfigure}

\begin{align}
\dfrac{dy_1}{dt} &= y_2 \\
\dfrac{dy_2}{dt} &= \mu(1-y_1^2)y_2 + y_1
\end{align}

The update steps for the Explicit Euler's method are simply
\begin{align}
	y_{1,i+1} &= 	y_{1,i} + \dfrac{dy_1}{dt}|_{i}h\\
	y_{2,i+1} &= 	y_{2,i} + \dfrac{dy_2}{dt}|_{i}h	
\end{align}

The update steps for the Implicit Euler's method are obtained by adding the derivative at the future step. These are given by
\begin{align}
y_{1,i+1} &= 	y_{1,i} + \dfrac{dy_1}{dt}|_{i+1}h\\
y_{2,i+1} &= 	y_{2,i} + \dfrac{dy_2}{dt}|_{i+1}h	
\end{align}

The above expressions are simplified as
\begin{align}
y_{1,i+1}  &= 	y_{1,i} + y_{2,i+1}h\\
y_{2,i+1}  &= 	y_{2,i} + (\mu(1-y_{1,i+1}^2)y_{2,i+1} + y_{1,i+1})h
\end{align}

From the above two equations, the updates are obtained by solving for $y_{1,i+1}$ and $y_{2,i+1}$. Since, these are a pair of non-linear equations, the solution is obtained using the \texttt{fsolve} function in Physika. These solutions are used as the update steps in Implicit Euler's method.
% y1 = zeros(n,1), y2 = zeros(n,1)
% y1[0] = 1, y2[0] = 1

\begin{lstlisting}[language=python]
def root2d(y, y1, y2, $\mu$, h):
 y1dash = y[0], y2dash = y[1]
 F[0] = -y1dash + y1 + (y2dash)*h
 F[1] = -y2dash + y2 + (($\mu$*(1-y1dash$^2$)*y2dash - y1dash))*h
 return F
\end{lstlisting}	

%# Explicit Euler's Method
%for i in range(2,n):
%  y1[i] = y1[i-1] + y2[i-1]*dt
%  y2[i] = y2[i-1] + ($\mu$*(1-y1[i-1]^2)*y2[i-1]-y1[i-1])*dt
\begin{lstlisting}[language=python]
def implicit_euler_method($\mu$: $\mathbb{R}$, t$_f$: $\mathbb{R}$, dt: $\mathbb{R}$):
  t = 0:dt:t$_f$
  n = len(t)
  y1i = zeros(n,1), y2i = zeros(n,1)
  y1i[0] = 1, y2i[0] = 1
  for i in 1 .. n-1:
    solution = fsolve($\lambda$ (y) root2d(y,y1i[i],y2i[i],$\mu$,dt),  [y1i[i], y2i[i]])
    y1i[i+1] = solution[0]
    y2i[i+1] = solution[1]     
\end{lstlisting}

Since the van der Pol equation with $\mu=5$ is a stiff system, the explicit method clearly fails in capturing the response of the variable. However, the implicit method still captures the oscillations in the response.
%\end{document}

%\section{Physika functions for stiff systems}
%MATLAB has a number of built-in functions for solving stiff systems of ODEs.

%\begin{marginfigure} 
%	\centering
%	\includegraphics[width = 2.5in]{figures/part1b/implicit_methods/ode23s.png}
%	\centering
%	\caption{Solution of the van der Pol equation using MATLAB function \texttt{ode23s}}
%	\label{fig:6}
%\end{marginfigure}

%\begin{marginfigure}
%\begin{verbatim}
%
%Syntax:
%[t,y] = ode15s(odefun,tspan,y0)
%[t,y] = ode23s(odefun,tspan,y0)
%[t,y] = ode23t(odefun,tspan,y0)
%[t,y] = ode23tb(odefun,tspan,y0)
%\end{verbatim}	
% 
%
%\texttt{odefun} is the function handle to the ODE to be solved
%
%\texttt{tspan} is the interval of integration %
%
%\texttt{y0} are the initial conditions
%
%
%\end{marginfigure}

%\begin{itemize}
%
%\item \texttt{ode15s} - This function is a variable-order solver based on numerical differentiation formulas. It is a multi-step solver that optionally uses the Gear backward differentiation formulas. This is used for stiff problems of low to medium accuracy.
%\item \texttt{ode23s}. This is based on a modified Rosenbrock formula of order two. As this is a single-step solver, it may be more efficient than \texttt{ode15s} at solving problems that permit crude tolerances or problems with solutions that change rapidly. It can solve some kinds of stiff problems for which \texttt{ode15s} is not effective. The \texttt{ode23s} solver evaluates the Jacobian during each step of the integration, so supplying it with the Jacobian matrix is critical to its reliability and efficiency.
%\item \texttt{ode23t}. This is an implementation of the trapezoidal rule using a "free" interpolant. This solver is preferred over \texttt{ode15s} if the problem is only moderately stiff and you need a solution without numerical damping. 
%\item \texttt{ode23tb}. This is an implementation of implicit Runge-Kutta formula with a trapezoidal rule step as its first stage and a backward differentiation formula of order two as its second stage. By construction, the same iteration matrix is used in evaluating both stages. Like \texttt{ode23s} and \texttt{ode23t}, this solver may be more efficient than \texttt{ode15s} for problems with crude tolerances.
%\end{itemize}

\end{document}